% ============================================================
%  DISCOVERY CURRICULUM — Journey 2, Module 5
%  "Differential Equations"
%  Compile with: pdflatex J2_M5.tex  (run twice for references)
% ============================================================

\documentclass[12pt, a4paper]{article}

\usepackage[a4paper, top=2.5cm, bottom=2.5cm, left=2.8cm, right=2.8cm, headheight=15pt]{geometry}
\usepackage{amsmath, amssymb, amsthm}
\usepackage{xcolor}
\usepackage{mdframed}
\usepackage{enumitem}
\usepackage{tikz}
\usetikzlibrary{arrows.meta, positioning, calc}
\usepackage{pgfplots}
\pgfplotsset{compat=1.18}
\usepackage{booktabs}
\usepackage{fancyhdr}
\usepackage{titlesec}
\usepackage{parskip}
\usepackage{hyperref}
\hypersetup{colorlinks=true, linkcolor=NavyBlue, urlcolor=NavyBlue}

% ── Colours ─────────────────────────────────────────────────
\definecolor{NavyBlue}{HTML}{1A3A5C}
\definecolor{MidBlue}{HTML}{2E6DA4}
\definecolor{Gold}{HTML}{C8922A}
\definecolor{LightBlue}{HTML}{EBF2FA}
\definecolor{LightGold}{HTML}{FDF4E3}
\definecolor{DarkGrey}{HTML}{3A3A3A}
\definecolor{BoxGreen}{HTML}{E8F5E9}
\definecolor{SidebarGrey}{HTML}{F5F5F5}

% ── Section formatting ───────────────────────────────────────
\titleformat{\section}{\Large\bfseries\color{NavyBlue}}{}{0em}{}[{\color{MidBlue}\vspace{2pt}\rule{\linewidth}{1.5pt}}]
\titleformat{\subsection}{\large\bfseries\color{MidBlue}}{}{0em}{}

% ── Header / Footer ─────────────────────────────────────────
\pagestyle{fancy}
\fancyhf{}
\fancyhead[L]{\small\color{DarkGrey}\textit{Discovery Curriculum}}
\fancyhead[R]{\small\color{DarkGrey}\textit{Journey 2 $\cdot$ Module 5}}
\fancyfoot[C]{\small\color{MidBlue}\thepage}
\renewcommand{\headrulewidth}{0.4pt}

% ── Environments ─────────────────────────────────────────────
\newmdenv[
  backgroundcolor=NavyBlue!7,
  linecolor=NavyBlue, linewidth=2pt,
  leftline=true, rightline=false, topline=false, bottomline=false,
  innerleftmargin=14pt, innerrightmargin=10pt,
  innertopmargin=10pt, innerbottommargin=10pt,
  skipabove=14pt, skipbelow=10pt
]{puzzlebox}

\newmdenv[
  backgroundcolor=LightGold,
  linecolor=Gold, linewidth=2pt,
  leftline=true, rightline=false, topline=false, bottomline=false,
  innerleftmargin=14pt, innerrightmargin=10pt,
  innertopmargin=10pt, innerbottommargin=10pt,
  skipabove=14pt, skipbelow=10pt
]{discoverybox}

\newmdenv[
  backgroundcolor=LightBlue,
  linecolor=MidBlue, linewidth=1.5pt,
  roundcorner=4pt,
  innerleftmargin=14pt, innerrightmargin=10pt,
  innertopmargin=10pt, innerbottommargin=10pt,
  skipabove=14pt, skipbelow=10pt
]{namebox}

\newmdenv[
  backgroundcolor=BoxGreen,
  linecolor=DarkGrey!35, linewidth=1pt,
  roundcorner=4pt,
  innerleftmargin=14pt, innerrightmargin=10pt,
  innertopmargin=10pt, innerbottommargin=10pt,
  skipabove=14pt, skipbelow=10pt
]{reflectionbox}

\newmdenv[
  backgroundcolor=SidebarGrey,
  linecolor=Gold, linewidth=1.2pt,
  leftline=true, rightline=false, topline=false, bottomline=false,
  innerleftmargin=12pt, innerrightmargin=10pt,
  innertopmargin=7pt, innerbottommargin=7pt,
  skipabove=10pt, skipbelow=10pt
]{thought}

\newcounter{problem}
\newenvironment{problem}[1][]{%
  \refstepcounter{problem}%
  \par\vspace{6pt}%
  \noindent\textbf{\color{MidBlue}Problem \theproblem}%
  \ifx&#1&\else\ \textit{(#1)}\fi\textbf{.}\enskip
}{\par\vspace{4pt}}

\newcounter{deepproblem}
\newenvironment{deepproblem}[1][]{%
  \refstepcounter{deepproblem}%
  \par\vspace{6pt}%
  \noindent\textbf{\color{NavyBlue}$\bigstar$\ Deep Problem \thedeepproblem}%
  \ifx&#1&\else\ \textit{(#1)}\fi\textbf{.}\enskip
}{\par\vspace{4pt}}

% ============================================================
\begin{document}
% ============================================================

\thispagestyle{empty}
\begin{center}
{\small\color{MidBlue}\textbf{DISCOVERY CURRICULUM}\quad$\cdot$\quad\textbf{JOURNEY 2: CHANGE}}

\vspace{0.6cm}
{\color{NavyBlue}\rule{\linewidth}{2pt}}
\vspace{0.5cm}

{\Huge\bfseries\color{NavyBlue}Module 5}\\[0.3cm]
{\LARGE\bfseries\color{MidBlue}Differential Equations}\\[0.4cm]
{\color{Gold}\rule{0.4\linewidth}{1.2pt}}\\[0.4cm]
{\large\itshape\color{DarkGrey}
Physics does not tell you where things are.\\[0.15cm]
It tells you how fast they are changing. The rest is mathematics.}

\vspace{1cm}
\begin{tabular}{rl}
\color{DarkGrey}\textbf{Journey}           & Journey 2 --- Change \\[4pt]
\color{DarkGrey}\textbf{Module}            & 5 of 6 \\[4pt]
\color{DarkGrey}\textbf{Prerequisites}     & J2 Modules 1--4 \\[4pt]
\color{DarkGrey}\textbf{You will discover} & ODEs, separation of variables, linear first-order, second-order, harmonic oscillator \\[4pt]
\color{DarkGrey}\textbf{Format}            & 8 problems $+$ 3 deep problems $\approx$ 14 pages \\
\end{tabular}

\vspace{1cm}
{\color{NavyBlue}\rule{\linewidth}{2pt}}
\end{center}

\vspace{0.5cm}

\noindent\textbf{The language of physics.}
Newton's laws do not say ``the ball is here.'' They say: ``the rate of change of momentum equals the applied force.'' That is a differential equation. Maxwell's equations do not describe the electric and magnetic fields directly --- they describe how those fields change in space and time. Einstein's field equations describe how spacetime curvature changes in response to matter. Schr\"{o}dinger's equation describes how a quantum state changes in time.

\medskip
In each case, the law is a differential equation. Solving it --- finding the function that satisfies it --- is the act of extracting physics from mathematics.

\newpage

% ============================================================
\section{Part 1 \texorpdfstring{---}{--} The Puzzle}
% ============================================================

\begin{puzzlebox}
\textbf{The Opening Puzzle.}

A cup of tea at $90°\text{C}$ is placed in a room at $20°\text{C}$. Newton's Law of Cooling states that the rate at which the tea cools is proportional to the difference between the tea's temperature and the room temperature.

\medskip
\textbf{Question 1.} Write this as a differential equation for $T(t)$, the temperature of the tea at time $t$. Use $k > 0$ for the proportionality constant.

\textbf{Question 2.} What is $T(0)$? What should $T(t)$ approach as $t \to \infty$? Does your equation produce this behaviour?

\textbf{Question 3.} If the tea cools from $90°\text{C}$ to $70°\text{C}$ in 5 minutes, find $k$. When will it reach $40°\text{C}$?

\textbf{Question 4.} Sketch $T(t)$ for $t \geq 0$. What kind of function does it look like?
\end{puzzlebox}

\begin{thought}
\textbf{What a differential equation is.}
A differential equation (DE) is an equation relating a function to its derivatives. It does not give you the function directly --- it gives you a rule that the function must obey at every point. Solving a DE means finding all functions that satisfy the rule. A DE usually has infinitely many solutions (a whole family), and a specific solution is selected by specifying initial conditions --- the value of the function (and possibly its derivatives) at one point. This is why physics requires both the law (DE) and initial conditions (state at $t=0$) to predict the future. The law constrains the family; the initial state selects the individual.
\end{thought}

% ============================================================
\section{Part 2 \texorpdfstring{---}{--} Exploration}
% ============================================================

\subsection{Step 1: Separation of variables}

\begin{problem}
\label{prob:separation}
The simplest class of ODEs separates: all $y$ terms on one side, all $x$ terms on the other.
\begin{enumerate}[label=(\alph*), topsep=2pt]
  \item Solve Newton's cooling equation $\frac{dT}{dt} = -k(T - T_{\text{room}})$ by substituting $u = T - T_{\text{room}}$, separating variables, and integrating both sides.
  \item Solve $\frac{dy}{dx} = y^2$ with $y(0) = 1$. At what finite time does the solution ``blow up''? What does this mean physically for a system governed by this equation?
  \item Solve the logistic equation $\frac{dP}{dt} = rP\!\left(1 - \frac{P}{K}\right)$ by separating variables and using partial fractions. Verify that the solution is $P(t) = \frac{K}{1 + Ae^{-rt}}$, the logistic function from Journey~1.
  \item Solve $\frac{dy}{dx} = \frac{x}{y}$. What curve does the solution describe geometrically?
  \item For which ODEs does separation of variables work? What structural property must the equation have?
\end{enumerate}
\end{problem}

\subsection{Step 2: Linear first-order ODEs}

\begin{problem}
\label{prob:linear_first}
A linear first-order ODE has the form $y' + P(x)y = Q(x)$.
\begin{enumerate}[label=(\alph*), topsep=2pt]
  \item If $Q(x) = 0$ (the homogeneous case), the equation is separable. Solve $y' + P(x)y = 0$.
  \item For the general case, multiply both sides by the \textbf{integrating factor} $\mu(x) = e^{\int P(x)\,dx}$. Show that the left side becomes $\frac{d}{dx}[\mu(x) y]$.
  \item Use this to solve $y' + \frac{1}{x}y = x$ for $x > 0$.
  \item A tank holds 100 litres of brine with 5 kg of salt dissolved. Fresh water flows in at 2 L/min and the mixture flows out at 2 L/min. Write and solve the DE for the amount of salt $S(t)$ at time $t$.
  \item What happens to $S(t)$ as $t \to \infty$? Does the answer make physical sense?
\end{enumerate}
\end{problem}

\begin{thought}
\textbf{The integrating factor idea.}
The integrating factor $\mu(x) = e^{\int P\,dx}$ appears to come from nowhere. But it is discoverable: you want to convert $y' + Py$ into a perfect derivative $(\mu y)'$. Expanding $(\mu y)' = \mu' y + \mu y'$, you need $\mu' = \mu P$, which gives $\mu = e^{\int P\,dx}$. The method designs its own solution. This is a recurring pattern in mathematics: when a problem seems stuck, ask what structure would make it tractable, then construct that structure. The integrating factor is an instance of the general technique of \emph{reducing to a form you can solve}.
\end{thought}

\subsection{Step 3: Second-order linear ODEs}

\begin{problem}
\label{prob:second_order}
The equation $y'' + by' + cy = 0$ with constants $b, c$ is a second-order linear homogeneous ODE.

\begin{enumerate}[label=(\alph*), topsep=2pt]
  \item Try a solution $y = e^{rx}$. Substitute and show that $r$ must satisfy the \textbf{characteristic equation} $r^2 + br + c = 0$.
  \item If the characteristic equation has two distinct real roots $r_1, r_2$: show that $y = Ae^{r_1 x} + Be^{r_2 x}$ is the general solution. How many free constants are there, and why?
  \item If it has a repeated real root $r$: the general solution is $y = (A + Bx)e^{rx}$. Verify this is a solution.
  \item If it has complex roots $r = \alpha \pm \beta i$: use Euler's formula $e^{i\theta} = \cos\theta + i\sin\theta$ to write the general solution in terms of real functions.
  \item Solve $y'' + 4y = 0$, $y(0) = 1$, $y'(0) = 0$. What function does this produce? Recognise it.
\end{enumerate}
\end{problem}

\begin{thought}
\textbf{Characteristic equations and linear recurrences.}
In Module~4 of Journey~1, you solved the Fibonacci recurrence by finding roots of a characteristic equation $r^2 = r + 1$. The method for second-order ODEs is identical in structure: assume an exponential solution, substitute, and obtain a polynomial whose roots determine everything. This is not a coincidence. Differential equations are the continuous-time analogues of recurrence relations. Every technique for one has an analogue for the other. The connection between sequences (discrete) and functions (continuous) runs through all of mathematics: from difference equations to differential equations, from $z$-transforms to Laplace transforms, from combinatorics to analysis.
\end{thought}

\subsection{Step 4: The harmonic oscillator}

\begin{problem}
\label{prob:harmonic}
A mass $m$ on a spring with spring constant $k$ satisfies \textbf{Newton's second law} $F = ma$:
\[
  m\ddot{x} = -kx \qquad \Longrightarrow \qquad \ddot{x} + \omega^2 x = 0, \quad \omega = \sqrt{k/m}.
\]
\begin{enumerate}[label=(\alph*), topsep=2pt]
  \item Solve the ODE. The general solution has two free constants --- determine them from $x(0) = A$ (initial displacement) and $\dot{x}(0) = 0$ (released from rest).
  \item The solution is $x(t) = A\cos(\omega t)$. Compute $\dot{x}(t)$ and $\ddot{x}(t)$. Show that the total energy $E = \frac{1}{2}m\dot{x}^2 + \frac{1}{2}kx^2$ is constant in time.
  \item Now add damping: $m\ddot{x} + b\dot{x} + kx = 0$ (friction proportional to velocity). Find the three qualitatively different types of motion depending on whether $b^2 - 4mk$ is positive, zero, or negative. Name them.
  \item Add a periodic driving force: $m\ddot{x} + kx = F_0\cos(\Omega t)$. If $\Omega \neq \omega$, the solution is a sum of two oscillations. If $\Omega = \omega$ (resonance), the solution grows without bound. Describe the resonance phenomenon and give a real-world example of its destructive consequences.
  \item The total energy in undamped free oscillation is constant. In damped oscillation, energy decreases. At what rate does energy decrease in the underdamped case? What happens to the energy?
\end{enumerate}
\end{problem}

\begin{center}
\begin{tikzpicture}
\begin{axis}[
  width=13cm, height=5.5cm,
  axis lines=center,
  xlabel={$t$}, ylabel={$x(t)$},
  xmin=0, xmax=14, ymin=-1.3, ymax=1.3,
  xtick={0,3.14,6.28,9.42,12.57},
  xticklabels={$0$,$\pi/\omega$,$2\pi/\omega$,$3\pi/\omega$,$4\pi/\omega$},
  ytick={-1,1},
  tick label style={font=\small, color=DarkGrey},
  grid=major, grid style={DarkGrey!12},
  title={\small\color{NavyBlue}Undamped ($\gamma=0$), underdamped ($\gamma=0.15$), overdamped ($\gamma=0.8$)},
  legend pos=north east,
  legend style={font=\small},
]
% Undamped
\addplot[NavyBlue, very thick, domain=0:14, samples=200] {cos(deg(x))};
% Underdamped
\addplot[Gold, thick, domain=0:14, samples=200] {e^(-0.15*x)*cos(deg(x))};
% Overdamped (no oscillation, just decay)
\addplot[MidBlue!60, thick, dashed, domain=0:14, samples=100] {e^(-0.8*x)*(cosh(deg(0.6*x))+1.33*sinh(deg(0.6*x)))};
\addlegendentry{undamped}
\addlegendentry{underdamped}
\addlegendentry{overdamped}
\end{axis}
\end{tikzpicture}
\end{center}

\begin{problem}
\label{prob:systems}
Many physical systems are described by \emph{systems} of coupled ODEs. Consider the predator-prey model (Lotka-Volterra):
\[
  \frac{dR}{dt} = aR - bRF, \qquad \frac{dF}{dt} = cRF - dF,
\]
where $R(t)$ is the rabbit population and $F(t)$ the fox population.
\begin{enumerate}[label=(\alph*), topsep=2pt]
  \item Interpret each term: what does $aR$ represent? What does $-bRF$ represent? What does $-dF$ represent?
  \item Find the equilibrium points (where $dR/dt = dF/dt = 0$). There are two: one where both populations are zero, and one where both are positive.
  \item The positive equilibrium is $R^* = d/c$, $F^* = a/b$. Explain why the equilibrium fox population depends on rabbit birth rate and vice versa.
  \item Near the positive equilibrium, the populations oscillate. Describe qualitatively what you expect from a phase portrait (plot of $R$ vs $F$). Do the oscillations match what we observe in real predator-prey systems?
  \item List three assumptions in the Lotka-Volterra model that are biologically unrealistic. For each, write one equation or term you would add to correct it.
\end{enumerate}
\end{problem}

\begin{problem}
\label{prob:slope_fields}
A slope field (direction field) gives a visual representation of the ODE $y' = f(x,y)$ by drawing a small arrow of slope $f(x,y)$ at each point $(x,y)$.
\begin{enumerate}[label=(\alph*), topsep=2pt]
  \item Sketch the slope field for $y' = y$ by computing the slope at a grid of points: $(\pm 1, \pm 1), (\pm 1, 0), (0, \pm 1)$, etc.
  \item On the same diagram, sketch a few solution curves. They should be tangent to the slope field everywhere.
  \item Sketch the slope field for $y' = -x/y$. What shape are the solution curves?
  \item \textbf{Euler's method}: given $y' = f(x,y)$ and initial condition $y(x_0) = y_0$, approximate the solution by taking small steps: $y_{n+1} = y_n + h\cdot f(x_n, y_n)$. Apply this to $y' = y$, $y(0) = 1$, with step size $h = 0.1$, for 10 steps. Compare with the exact solution.
\end{enumerate}
\end{problem}

% ============================================================
\section{Part 3 \texorpdfstring{---}{--} Discovery}
% ============================================================

\begin{discoverybox}
\textbf{What have we built?}

\textbf{First.} A differential equation is a constraint on a function: it specifies not the function itself but the relationship between the function and its rates of change. Solutions form a family; initial conditions select a specific member.

\textbf{Second.} Separation of variables, integrating factors, and the characteristic equation method are the three main tools for first- and second-order linear ODEs. Each works because it reduces the problem to integration.

\textbf{Third.} The harmonic oscillator $\ddot{x} + \omega^2 x = 0$ is the central equation of physics. Its solutions $\cos(\omega t)$ and $\sin(\omega t)$ describe everything that oscillates: springs, pendulums, electric circuits, sound waves, light, and quantum mechanical states.

\textbf{Fourth.} The characteristic equation for a second-order ODE is structurally identical to the characteristic equation for a second-order recurrence. The exponential function $e^{rx}$ is the fundamental building block of both.

\medskip
\textbf{The deep question.} You have found solutions to linear ODEs. What makes linearity so important? The key: if $y_1$ and $y_2$ are solutions to a \emph{linear} ODE, then $Ay_1 + By_2$ is also a solution. This is the \emph{superposition principle}. It is the reason that the solutions form a vector space (Journey~3 will make this precise). Does the superposition principle hold for the logistic equation? What does that tell you about the difficulty of nonlinear ODEs?
\end{discoverybox}

% ============================================================
\section{Part 4 \texorpdfstring{---}{--} Naming the Ideas}
% ============================================================

\begin{namebox}
\textbf{Ordinary differential equations.}

An \textbf{ODE} is an equation involving a function $y(x)$ and its derivatives. The \textbf{order} is the highest derivative appearing. A \textbf{linear} ODE has the form $a_n y^{(n)} + \cdots + a_1 y' + a_0 y = f(x)$; \textbf{nonlinear} otherwise.

\medskip
\textbf{First-order separable}: $\frac{dy}{dx} = g(x)h(y)$ $\Rightarrow$ $\int\frac{dy}{h(y)} = \int g(x)\,dx$.

\textbf{First-order linear}: $y' + P(x)y = Q(x)$ $\Rightarrow$ multiply by $\mu = e^{\int P\,dx}$.

\textbf{Second-order linear constant-coefficient}: $y'' + by' + cy = 0$.\\
Characteristic equation $r^2 + br + c = 0$ with roots $r_{1,2}$.\\
--- Two real roots: $y = Ae^{r_1 x} + Be^{r_2 x}$.\\
--- Repeated root $r$: $y = (A+Bx)e^{rx}$.\\
--- Complex roots $\alpha\pm\beta i$: $y = e^{\alpha x}(A\cos\beta x + B\sin\beta x)$.
\end{namebox}

\begin{namebox}
\textbf{The harmonic oscillator.}

$\ddot{x} + \omega^2 x = 0$ has solutions $x(t) = A\cos(\omega t) + B\sin(\omega t)$.

With damping $\ddot{x} + 2\gamma\dot{x} + \omega^2 x = 0$:
\begin{itemize}[noitemsep, topsep=4pt]
  \item $\gamma < \omega$: \textbf{underdamped} --- decaying oscillation
  \item $\gamma = \omega$: \textbf{critically damped} --- fastest non-oscillatory decay
  \item $\gamma > \omega$: \textbf{overdamped} --- slow exponential decay
\end{itemize}
\textbf{Resonance} occurs when driving frequency equals natural frequency $\omega$; amplitude grows without bound in an undamped system.
\end{namebox}

% ============================================================
\section{Part 5 \texorpdfstring{---}{--} Deep Problems}
% ============================================================

\begin{deepproblem}[Existence and uniqueness]
The \textbf{Picard-Lindel\"{o}f theorem} guarantees that if $f(x,y)$ is continuous and Lipschitz in $y$ near $(x_0, y_0)$, then the IVP $y' = f(x,y)$, $y(x_0) = y_0$ has a unique solution near $x_0$.
\begin{enumerate}[label=(\alph*)]
  \item The ODE $y' = y^{2/3}$, $y(0) = 0$ has two solutions: $y = 0$ and $y = (x/3)^3$. Find both. Why does this not contradict Picard-Lindel\"{o}f?
  \item The solution to $y' = y^2$, $y(0) = 1$ blows up in finite time. Does this contradict the theorem?
  \item Picard iteration: define $y_0(x) = y_0$ (constant), then $y_{n+1}(x) = y_0 + \int_{x_0}^x f(t, y_n(t))\,dt$. Apply this to $y' = y$, $y(0) = 1$ for 4 iterations. What sequence of functions do you obtain?
  \item What does the Picard iteration for $y' = y$, starting from $y_0 = 1$, produce in the limit? Connect to the series for $e^x$.
\end{enumerate}
\end{deepproblem}

\begin{deepproblem}[The heat equation as a limit]
The \textbf{heat equation} $\frac{\partial u}{\partial t} = \alpha\frac{\partial^2 u}{\partial x^2}$ governs temperature $u(x,t)$ along a rod.
\begin{enumerate}[label=(\alph*)]
  \item This is a \emph{partial} differential equation (PDE): the unknown function depends on both $x$ and $t$. Try a \textbf{separated} solution $u(x,t) = X(x)T(t)$. Substitute and divide by $XT$ to show that $T'/\alpha T = X''/X = -\lambda$ (a constant).
  \item Solve the two resulting ODEs. What type of function is $T(t)$? What type is $X(x)$?
  \item For a rod of length $L$ with $u(0,t) = u(L,t) = 0$ (ends held at zero temperature), the boundary conditions force $X(x) = \sin(n\pi x/L)$ for integer $n$. Write the general separated solution.
  \item The full solution is a superposition $u(x,t) = \sum_{n=1}^\infty B_n \sin\!\left(\frac{n\pi x}{L}\right) e^{-\alpha(n\pi/L)^2 t}$. For large $t$, which terms dominate? What does the temperature distribution approach?
  \item The coefficients $B_n$ are determined by the initial condition $u(x,0) = f(x)$ via $B_n = \frac{2}{L}\int_0^L f(x)\sin(n\pi x/L)\,dx$. This is a \textbf{Fourier sine series}. The idea that \emph{any} initial distribution can be written as a sum of sines is the starting point of Journey~6. Does it surprise you that this is possible?
\end{enumerate}
\end{deepproblem}

\begin{deepproblem}[Chaos in three coupled ODEs: the Lorenz system]
The Lorenz system is three coupled first-order ODEs:
\[
  \dot{x} = \sigma(y - x), \qquad
  \dot{y} = x(\rho - z) - y, \qquad
  \dot{z} = xy - \beta z,
\]
with parameters $\sigma = 10$, $\rho = 28$, $\beta = 8/3$.
\begin{enumerate}[label=(\alph*)]
  \item Find the equilibria (fixed points) of the Lorenz system by setting all three derivatives to zero. How many are there?
  \item The system conserves no energy and is dissipative (volumes in phase space shrink). Yet for $\rho > 1$, solutions do not converge to a fixed point. What else could they do?
  \item Lorenz discovered numerically that two solutions starting from almost identical initial conditions diverge exponentially. Starting from $(1, 1, 1)$ vs.\ $(1 + 10^{-5}, 1, 1)$: after 30 time units, the solutions look completely different. What does this imply about weather prediction?
  \item The set of all long-run positions forms the \textbf{Lorenz attractor} --- a set of fractal dimension $\approx 2.06$, neither a surface nor a solid. Describe what ``fractal dimension'' means informally.
  \item \textit{(Research question.)} Chaos in the Lorenz system is governed by the \textbf{Lyapunov exponent} $\lambda > 0$: nearby trajectories diverge as $e^{\lambda t}$. This sets a fundamental limit on predictability: errors double every $\sim \ln 2/\lambda$ time units. For weather, this is $\approx 5$ days. Does this mean long-range weather forecasting is fundamentally impossible, or merely practically difficult?
\end{enumerate}
\end{deepproblem}

% ============================================================
\section{Part 6 \texorpdfstring{---}{--} Reflection}
% ============================================================

\begin{reflectionbox}
\textbf{Your reflection task.}

\begin{enumerate}
  \item \textbf{On physics as differential equations.} You have now seen Newton's second law, Newton's cooling, logistic population growth, and the harmonic oscillator --- all as ODEs. In each case, the physical law states a \emph{local} constraint (what is happening right now, at this instant), and the solution describes the \emph{global} behaviour (what happens over all time). Write a paragraph: is this decomposition of physics into local law plus global solution illuminating, or does it obscure something about how the world works?

  \item \textbf{On the superposition principle.} Linear ODEs obey superposition: sums of solutions are solutions. Nonlinear ODEs do not. The logistic equation, the Lorenz system, and the equations governing fluid turbulence are all nonlinear --- and all much harder. Is linearity an intrinsic property of certain physical systems, or is it an approximation we make because nonlinear systems are too hard to analyse?

  \item \textbf{On chaos.} The Lorenz system is deterministic: the future is completely determined by the present. Yet it is unpredictable in practice. Does this distinction --- deterministic but unpredictable --- feel meaningful or hollow to you? Does it change your view of free will, or of the possibility of a ``theory of everything''?

  \item \textbf{Looking ahead.} Module~6 opens the door to calculus in higher dimensions: partial derivatives, gradients, and the geometry of change in space. The heat equation in the deep problems is an example of a PDE --- a differential equation where the unknown depends on multiple variables. What new difficulties do you anticipate when calculus moves into higher dimensions?
\end{enumerate}
\end{reflectionbox}

\medskip
\noindent\textbf{Looking ahead.}
Module~6 closes Journey~2 by extending calculus to functions of several variables. The derivative becomes a vector (the gradient). The integral becomes a double or triple integral over a region. And new phenomena appear that have no single-variable analogue --- curl, divergence, and the theorems of Green, Stokes, and Gauss --- which will be the foundation of Journey~3's study of fields and space.

\vspace{1.4cm}
\begin{center}
{\color{MidBlue}\rule{0.5\linewidth}{0.8pt}}\\[6pt]
{\small\color{DarkGrey}\textit{End of Journey 2, Module 5 --- Differential Equations}}\\[4pt]
{\small\color{DarkGrey}Continue to: \textbf{Module 6 --- Multivariable Calculus}}\\[6pt]
{\color{MidBlue}\rule{0.5\linewidth}{0.8pt}}
\end{center}

\end{document}